\section*{Background}

Graph classification on proteins and molecules has become a standard use case for graph neural networks because graph structure preserves both local biochemical motifs and broader relational organization. In biomolecular settings, this dual dependence matters: a graph label may reflect coordinated substructures rather than one dominant neighborhood pattern, so useful prediction signals can be distributed across depth, across branches, and across graph-level summaries. The background below therefore moves from the application setting to the mechanism setting that motivates CR-GNN.

\subsection*{Graph neural networks for biomolecular graph classification}

Graph-based learning is now routinely used for protein and molecular tasks in which the target depends on structured relational information rather than on independent tabular descriptors alone \cite{kipf2016semi,gilmer2017neural,hamilton2017inductive}. In biomolecular applications, graph neural networks have already been used for structure-aware protein function prediction, protein interface modeling, and related molecular learning problems where the relation among substructures is as important as the substructures themselves \cite{gligorijevic2021structure,you2021deepgraphgo,reau2022deeprankgnn}. Benchmark suites such as PROTEINS, DD, ENZYMES, MUTAG, AIDS, and Mutagenicity remain widely used because they expose differences in graph size, sparsity, and classification difficulty while retaining direct biological or molecular interpretation \cite{fey2019fast,morris2020tudataset}. In the present study, these datasets are used as a reproducible biomolecular graph-classification package rather than as a proxy for downstream biological discovery.

\subsection*{Residual reuse and deep message passing}

Deeper message passing can enlarge receptive fields, but it also increases the risk of representation degradation. Over-smoothing and related depth effects can suppress informative intermediate states, especially when repeated propagation progressively attenuates local variation \cite{li2018deeper,oono2020simple,chen2023revisiting}. Residual connections, normalization, and regularization alleviate this problem \cite{he2016deep,bresson2017residual,zhao2020pair,cai2021graphnorm,rong2019dropedge}, but most residual designs still reuse information within a single branch. As a result, they improve optimization without directly testing whether alternative information-flow topologies preserve more useful historical signals.

\subsection*{Cross-branch interaction and graph-level feedback}

Cross-layer reuse in graph learning has been explored through skip connections, dense aggregation, APPNP-style propagation, and Jumping Knowledge readout \cite{xu2018representation,klicpera2018combining,gao2019graph}. These approaches preserve information from multiple depths, yet they do not directly isolate whether exchanging historical states across branches or re-injecting graph summaries across block boundaries is preferable to ordinary same-path residual reuse. CR-GNN is designed around exactly that distinction. It treats node-level cross exchange and graph-level cross-summary propagation as explicit alternatives to standard residual carryover, so that the effect of information-flow topology can be compared under one benchmark protocol.

\subsection*{Residual routing as a mechanism question}

The present paper studies not only whether residual information is reused, but also how it is routed. A dense learnable residual gate is one possible design, but filtered routing may behave differently when the target architecture is cross-oriented or when the residual path is already strong. This motivates the residual-routing analysis in the later sections, where learnable, top-k, and sparse routing modes are evaluated under the same experimental protocol. In this sense, the paper extends a pure architecture comparison into a mechanism study of residual injection strength and selectivity.

\subsection*{Study positioning}

The paper is best understood as a biomolecular graph-classification methods study with a mechanism-oriented empirical focus. It does not claim universal superiority of cross-residual interaction, and it does not attempt direct biological discovery. Instead, it asks a narrower and more defensible question: when does cross-residual information flow help beyond plain stacking, when does ordinary residual reuse remain the better default, and how does residual routing strategy change that answer? That positioning is important for the interpretation of the results: the goal is to characterize the boundary of usefulness for cross-residual and routed-residual designs, not to force a single winner narrative.
